\documentclass[aspectratio=169,11pt]{beamer}

\usepackage[T1]{fontenc}
\usepackage[utf8]{inputenc}
\usepackage{booktabs}
\usepackage{amssymb}
\usepackage{array}
\usepackage{ragged2e}
\usepackage{helvet}
\renewcommand{\familydefault}{\sfdefault}

\makeatletter
\def\input@path{{docs/rub-beamer/}{rub-beamer/}}
\makeatother
\usetheme[empty]{Rub}

\colorlet{Ink}{saphierblau}
\definecolor{Muted}{HTML}{66717C}
\colorlet{Blue}{saphierblau}
\colorlet{BlueLight}{saphierblau!10}
\colorlet{Green}{gelbgruen!55!black}
\colorlet{GreenLight}{gelbgruen!15}
\definecolor{Amber}{HTML}{A46300}
\definecolor{AmberLight}{HTML}{FFF0D2}
\colorlet{Red}{alertred}
\colorlet{RedLight}{alertred!12}
\colorlet{GreyLight}{lichtgrau}

\setbeamercolor{progress done}{bg=saphierblau,fg=white}
\setbeamercolor{progress remaining}{bg=lichtgrau,fg=saphierblau}
\setbeamerfont{title}{series=\bfseries,size=\fontsize{24}{28}}
\setbeamerfont{subtitle}{series=\normalfont,size=\fontsize{13}{16}}
\setbeamerfont{frametitle}{series=\bfseries,size=\fontsize{16}{18}}
\setbeamerfont{normal text}{size=\fontsize{10.5}{13}}
\AtBeginEnvironment{column}{\RaggedRight}

\newcommand{\metric}[3]{%
  \begin{beamercolorbox}[rounded=true,sep=0.42em,wd=0.90\linewidth]{#3}
    {\fontsize{17}{19}\selectfont\bfseries #1}\par
    \vspace{0.04cm}{\footnotesize #2}
  \end{beamercolorbox}}

\newcommand{\done}{\colorbox{GreenLight}{\textcolor{Green}{\strut\bfseries complete}}}
\newcommand{\running}{\colorbox{BlueLight}{\textcolor{Blue}{\strut\bfseries running}}}
\newcommand{\validated}{\colorbox{GreenLight}{\textcolor{Green}{\strut\bfseries implemented}}}
\newcommand{\repair}{\colorbox{RedLight}{\textcolor{Red}{\strut\bfseries no-go}}}
\newcommand{\partialstatus}{\colorbox{AmberLight}{\textcolor{Amber}{\strut\bfseries partial}}}
\newcommand{\scaffold}{\colorbox{GreyLight}{\textcolor{Muted}{\strut\bfseries scaffold}}}

\title{Mitigating factual hallucinations}
\subtitle{What we are researching, why the study is designed this way, and where implementation stands}
\author{Master's research project}
\date{17 July 2026}

\begin{document}

\begin{frame}[plain]
  \vspace{0.55cm}
  {\color{Blue}\rule{1.5cm}{3pt}}\\[0.45cm]
  {\usebeamerfont{title}\inserttitle\par}
  \vspace{0.25cm}
  {\usebeamerfont{subtitle}\insertsubtitle\par}
  \vfill
  \begin{beamercolorbox}[rounded=true,sep=0.7em,wd=0.91\textwidth]{progress remaining}
    \textbf{Current position:} the E0--E10 software, provenance, replay, and analysis path is implemented. The earlier Bonsai pilot is preserved but superseded; live Qwen validation and every Qwen scientific phase remain to be executed.
  \end{beamercolorbox}
  \vspace{0.4cm}
  {\small\color{Muted}Evidence snapshot: 17 July 2026}
\end{frame}

\begin{frame}{Reliability means avoiding two opposite failures}
  \footnotesize
  \begin{columns}[T,onlytextwidth]
    \begin{column}{0.46\textwidth}
      \textcolor{Red}{\bfseries Confident fabrication}\par\vspace{0.15cm}
      The model gives a fluent factual answer that is wrong. The user may trust it precisely because it sounds authoritative.\par\vspace{0.25cm}
      \textbf{What we call a hallucination}\par
      An \emph{incorrect attempted answer}: the response makes a factual commitment that does not match the accepted answer.
    \end{column}
    \begin{column}{0.46\textwidth}
      \textcolor{Amber}{\bfseries Unnecessary refusal}\par\vspace{0.15cm}
      The model says ``I do not know'' even though it could answer correctly. Hallucinations fall, but usefulness falls with them.\par\vspace{0.25cm}
      \textbf{Why error rate alone misleads}\par
      A model that refuses every question has zero attempted hallucinations and zero practical utility.
    \end{column}
  \end{columns}
  \vspace{0.2cm}
  \begin{beamercolorbox}[rounded=true,sep=0.3em]{progress done}
    \textbf{Research goal:} answer when the model is likely correct, intervene when knowledge may be recoverable, and abstain only when residual risk remains too high.
  \end{beamercolorbox}
\end{frame}

\begin{frame}{We frame the task as selective prediction}
  \small
  \begin{columns}[T,onlytextwidth]
    \begin{column}{0.48\textwidth}
      \textbf{Optimization target}\par\vspace{0.18cm}
      Maximize answer coverage subject to a strict bound on incorrect attempted answers, while preserving utility and safety.\par\vspace{0.45cm}
      \begin{beamercolorbox}[rounded=true,sep=0.45em]{progress remaining}
        \centering
        $\max\ \mathrm{Coverage}$\quad subject to\quad
        $\mathrm{HallucinationRisk}\leq\epsilon$
      \end{beamercolorbox}
    \end{column}
    \begin{column}{0.46\textwidth}
      \textbf{Why this formulation}\par\vspace{0.18cm}
      \begin{itemize}\setlength\itemsep{0.12cm}
        \item It makes refusal cost visible.
        \item It supports fair comparisons at matched coverage.
        \item It separates ``fewer errors'' from ``more knowledge.''
        \item It allows a statistical upper bound when zero errors are observed.
      \end{itemize}
    \end{column}
  \end{columns}
  \vspace{0.3cm}
  {\footnotesize\color{Muted}Example: zero errors in 1,000 attempted answers implies an approximate one-sided 95\% upper risk bound of $3/1000=0.3\%$, not universal zero risk.}
\end{frame}

\begin{frame}{Four questions determine whether steering is genuinely useful}
  \small
  \begin{columns}[T,onlytextwidth]
    \begin{column}{0.48\textwidth}
      \textcolor{Blue}{\bfseries RQ1 --- Adaptive vs. static}\par
      Can the query choose its vector, layer, and strength better than one fixed direction?\par\vspace{0.35cm}
      \textcolor{Blue}{\bfseries RQ2 --- Knowledge vs. abstention}\par
      Do incorrect answers become correct, or merely become refusals?
    \end{column}
    \begin{column}{0.48\textwidth}
      \textcolor{Blue}{\bfseries RQ3 --- Sparse and protected edits}\par
      Can factuality features be changed without harming language, safety, or general capability?\par\vspace{0.35cm}
      \textcolor{Blue}{\bfseries RQ4 --- Prompt interaction}\par
      Does internal steering add value beyond inexpensive textual instructions?
    \end{column}
  \end{columns}
  \vspace{0.35cm}
  \begin{beamercolorbox}[rounded=true,sep=0.45em]{progress remaining}
    \textbf{Why all four are needed:} a lower raw hallucination count can otherwise be explained by stronger perturbations, more refusals, benchmark-specific tuning, or a prompt that already produces the same effect.
  \end{beamercolorbox}
\end{frame}

\begin{frame}{Steering changes behavior inside inference}
  \footnotesize
  \begin{columns}[T,onlytextwidth]
    \begin{column}{0.45\textwidth}
      \textbf{What it is}\par\vspace{0.15cm}
      We read an internal activation, add a learned direction at a chosen transformer site, and continue generation without retraining all model weights.\par\vspace{0.25cm}
      \textbf{Methods being compared}\par
      M0 no steering; M1 dense centroid; M2 published baselines; M3 adaptive controller; M4 sparse features; M5 protected subspace; M6 composite policy.
    \end{column}
    \begin{column}{0.49\textwidth}
      \textbf{Why go beyond one global vector}\par\vspace{0.15cm}
      A dense direction can mix factuality with uncertainty, refusal style, language, and safety. Applying it to every query may therefore fix some errors while creating new regressions.\par\vspace{0.25cm}
      \textbf{Planned improvement}\par
      Estimate risk from the current query, intervene only in the recoverable region, and constrain the edit to features that preserve protected behavior.
    \end{column}
  \end{columns}
  \vspace{0.18cm}
  \begin{beamercolorbox}[rounded=true,sep=0.3em]{progress done}
    \textbf{Core causal test:} compare every intervention with the same checkpoint, questions, prompts, decoding, and grader; change only the internal policy.
  \end{beamercolorbox}
\end{frame}

\begin{frame}{Outcome transitions reveal what an improvement means}
  \small
  \begin{columns}[T,onlytextwidth]
    \begin{column}{0.31\textwidth}
      \textbf{Response labels}\par\vspace{0.15cm}
      C --- correct\par
      P --- partially correct\par
      I --- incorrect attempted\par
      A --- abstention\par
      U --- unscorable
    \end{column}
    \begin{column}{0.31\textwidth}
      \textbf{Desired transitions}\par\vspace{0.15cm}
      $I\rightarrow C$: knowledge recovery\par
      $I\rightarrow A$: safer abstention\par
      $C\rightarrow C$: correct preservation
    \end{column}
    \begin{column}{0.31\textwidth}
      \textbf{Failure transitions}\par\vspace{0.15cm}
      $C\rightarrow A$: strict over-refusal\par
      $C\rightarrow I$: factual regression\par
      correct $\rightarrow$ wrong language or unsafe behavior
    \end{column}
  \end{columns}
  \vspace{0.35cm}
  \begin{beamercolorbox}[rounded=true,sep=0.45em]{progress remaining}
    \textbf{Why forced-answer and likelihood tests matter:} if gains disappear when abstention is discouraged and gold-answer likelihood does not rise, steering probably changed refusal behavior rather than knowledge expression.
  \end{beamercolorbox}
\end{frame}

\begin{frame}{Three benchmarks serve different scientific roles}
  \footnotesize
  \renewcommand{\arraystretch}{1.12}
  \begin{tabular}{@{}>{\bfseries\raggedright\arraybackslash}p{0.22\textwidth} >{\raggedright\arraybackslash}p{0.25\textwidth} >{\raggedright\arraybackslash}p{0.43\textwidth}@{}}
    \toprule
    Benchmark & Role & Why it is included \\
    \midrule
    TriviaQA & 5,000-question controller split & Large development source for probes, vectors, routers, and in-distribution testing; evidence documents are discarded to keep the task closed-book. \\
    SimpleQA Verified & 1,000 untouched questions & External short-answer test with an explicit correct/incorrect/not-attempted rubric; useful for detecting refusal-driven gains. \\
    AA-Omniscience Public-600 & 600 untouched questions & Directly tests calibrated knowledge boundaries and includes an official abstention-aware setup. Public size limits domain-level claims. \\
    \bottomrule
  \end{tabular}
  \vspace{0.16cm}
  \begin{beamercolorbox}[rounded=true,sep=0.25em]{progress done}
    \textbf{Why this split matters:} methods are learned on TriviaQA and transferred unchanged to the two evaluation sets, reducing target-benchmark overfitting.
  \end{beamercolorbox}
\end{frame}

\begin{frame}{Resource limits changed the execution platform}
  \footnotesize
  \begin{columns}[T,onlytextwidth]
    \begin{column}{0.47\textwidth}
      \textbf{What happened}\par\vspace{0.15cm}
      \begin{enumerate}\setlength\itemsep{0.06cm}
        \item Colab Free offered a T4 or v5e TPU, but its RAM/session limits could not support the complete activation workflow reliably.
        \item The Colab notebook and import bundle were retired before scientific E1 results.
        \item The study moved to one native Apple-silicon checkpoint on the available MacBook Pro.
      \end{enumerate}
    \end{column}
    \begin{column}{0.47\textwidth}
      \textbf{Why this was the defensible choice}\par\vspace{0.15cm}
      \begin{itemize}\setlength\itemsep{0.06cm}
        \item Stable access to internal hooks is more important than nominal accelerator speed.
        \item Local runs can resume without session eviction.
        \item One pinned representation prevents mixing incompatible activations across runtimes.
        \item The causal questions remain testable within one model.
      \end{itemize}
    \end{column}
  \end{columns}
  \vspace{0.15cm}
  {\scriptsize\color{Red}\textbf{Limitation:} conclusions apply only to this checkpoint/runtime; cross-model replication remains future work.}
\end{frame}

\begin{frame}{Uniform four-bit Qwen is the frozen scientific target}
  \small
  \begin{columns}[T,onlytextwidth]
    \begin{column}{0.48\textwidth}
      \textbf{Pinned system}\par\vspace{0.15cm}
      \texttt{mlx-community/Qwen3.6-27B-4bit}\par
      exact revision \texttt{c000ac2c...3282}\par
      official MLX 0.31.2 / MLX-LM 0.31.3\par
      Apple M4 Max, 48 GiB unified memory\par\vspace{0.35cm}
      64 blocks, hidden size 5,120; pinned 16.08 GB snapshot; uniform affine group-64 four-bit representation of \texttt{Qwen/Qwen3.6-27B}.
    \end{column}
    \begin{column}{0.46\textwidth}
      \textbf{Why this model}\par\vspace{0.15cm}
      \begin{itemize}\setlength\itemsep{0.13cm}
        \item It fits the approved 48 GiB machine with headroom for inference state.
        \item Native MLX uses the released Apple-silicon representation.
        \item One uniform quantization rule avoids layer-dependent precision as a hidden experimental factor.
        \item Live preflight must prove all intervention sites on both linear- and full-attention blocks before E0.
      \end{itemize}
    \end{column}
  \end{columns}
\end{frame}

\begin{frame}{Human review closes a gap automation cannot settle}
  \small
  \begin{columns}[T,onlytextwidth]
    \begin{column}{0.47\textwidth}
      \textbf{What we did}\par\vspace{0.15cm}
      Exact, fuzzy, and semantic matching created a blinded queue of the 200 highest-risk contamination candidates. A human reviewed the supplied CSV without model-condition information.\par\vspace{0.35cm}
      \textbf{Result}\par
      200 rows reviewed; 197 distinct candidate pairs; 3 genuine overlaps were excluded; none affected the 500-question E0 cohort.
    \end{column}
    \begin{column}{0.47\textwidth}
      \textbf{Why a person was needed}\par\vspace{0.15cm}
      Similar wording, shared entities, and paraphrases can look like duplicates to an embedding model without asking the same factual question. Conversely, true leakage can be phrased differently.\par\vspace{0.35cm}
      \textbf{Why it gates E1}\par
      Vectors and controllers must not learn from evaluation questions. The signed review receipt proves the development/evaluation boundary before expensive computation starts.
    \end{column}
  \end{columns}
\end{frame}

\begin{frame}{External graders are used where exact matching is insufficient}
  \small
  \begin{columns}[T,onlytextwidth]
    \begin{column}{0.47\textwidth}
      \textbf{What we implemented}\par\vspace{0.15cm}
      Local deterministic generation is followed by frozen OpenRouter grading for semantic benchmarks: GPT-4.1 with the released SimpleQA rubric and stable Gemini 2.5 Flash with the AA rubric.\par\vspace{0.35cm}
      The adapter records the prompt, model identity, transport, retries, parser result, and terminal failures.
    \end{column}
    \begin{column}{0.47\textwidth}
      \textbf{Why this is necessary}\par\vspace{0.15cm}
      Correct answers often use aliases or explanatory wording that string matching cannot judge. AA also needs a partial-credit distinction.\par\vspace{0.35cm}
      \textbf{Why it is frozen and fail-closed}\par
      A changing grader can change the experiment. Atomic labels, bounded retries, and ``unscorable'' on failure prevent silent defaults from biasing results.
    \end{column}
  \end{columns}
  \vspace{0.3cm}
  {\footnotesize\color{Muted}The retired Gemini preview endpoint was replaced before E1 grading; the amendment preserves the rubric but explicitly limits checkpoint comparability.}
\end{frame}

\begin{frame}{Content-addressed gates keep runs auditable}
  \footnotesize
  \begin{columns}[T,onlytextwidth]
    \begin{column}{0.46\textwidth}
      \textbf{What was built}\par\vspace{0.15cm}
      Every phase has frozen inputs, manifests, SHA-256 identities, append-only records, resumable chain heads, prerequisite receipts, and replayable finalizers.\par\vspace{0.25cm}
      Model, tokenizer, prompt template, decoding, hardware, runtime, split, grader, and code identities are bound to each run.
    \end{column}
    \begin{column}{0.48\textwidth}
      \textbf{Why this is unusually strict}\par\vspace{0.15cm}
      \begin{itemize}\setlength\itemsep{0.06cm}
        \item Multi-day local runs must resume without duplicating or silently changing rows.
        \item Later phases must not consume unverified upstream outputs.
        \item Hyperparameters must be selected before untouched tests are opened.
        \item A result should be independently replayable from immutable evidence.
      \end{itemize}
    \end{column}
  \end{columns}
  \vspace{0.15cm}
  \begin{beamercolorbox}[rounded=true,sep=0.28em]{progress done}
    \textbf{Interpretation rule:} passing tests shows that software enforces the protocol; only a completed scientific ledger supports an experimental claim.
  \end{beamercolorbox}
\end{frame}

\begin{frame}{E0 tests whether local experimentation is viable}
  \footnotesize
  \begin{columns}[T,onlytextwidth]
    \begin{column}{0.47\textwidth}
      \textbf{What E0 tested}\par\vspace{0.15cm}
      \begin{itemize}\setlength\itemsep{0.04cm}
        \item exact model and snapshot identity
        \item tokenizer and chat-template identity
        \item deterministic decoding across repeated questions
        \item throughput, latency, and peak unified memory
        \item all required activation capture/intervention sites
        \item scientific admission after contamination review
      \end{itemize}
    \end{column}
    \begin{column}{0.47\textwidth}
      \textbf{Why E0 came first}\par\vspace{0.15cm}
      The rest of the plan depends on stable activations. If the checkpoint, template, cached path, or hook changes, learned vectors no longer refer to the same internal system.\par\vspace{0.2cm}
      Measuring memory before activation collection also prevents designing E2--E8 around an impossible storage or compute budget.
    \end{column}
  \end{columns}
  \vspace{0.14cm}
  {\scriptsize\color{Muted}\textbf{Design choice:} 500 model-independent reserved questions, each generated twice with the exact E1 decoding path.}
\end{frame}

\begin{frame}{The Bonsai E0 pilot is preserved but superseded}
  \begin{columns}[T,onlytextwidth]
    \begin{column}{0.72\textwidth}
      \begin{columns}[T,onlytextwidth]
        \begin{column}{0.5\textwidth}\metric{500}{questions, each repeated twice}{progress remaining}\end{column}
        \begin{column}{0.5\textwidth}\metric{0}{determinism mismatches}{progress remaining}\end{column}
      \end{columns}
      \vspace{0.2cm}
      \begin{columns}[T,onlytextwidth]
        \begin{column}{0.5\textwidth}\metric{13.84}{mean generated tokens/s}{progress remaining}\end{column}
        \begin{column}{0.5\textwidth}\metric{4.86 GB}{measured peak memory}{progress remaining}\end{column}
      \end{columns}
    \end{column}
    \begin{column}{0.25\textwidth}
      \begin{beamercolorbox}[rounded=true,sep=0.4em]{progress remaining}
        \small\textbf{Pilot passed}\par\vspace{0.12cm}
        \footnotesize
        \textcolor{Green}{\checkmark}\ checkpoint\par
        \textcolor{Green}{\checkmark}\ chat template\par
        \textcolor{Green}{\checkmark}\ deterministic decode\par
        \textcolor{Green}{\checkmark}\ MLX runtime
      \end{beamercolorbox}
    \end{column}
  \end{columns}
  \vspace{0.08cm}
  \begin{beamercolorbox}[rounded=true,sep=0.32em,wd=0.95\textwidth]{progress done}
    \small\textbf{Meaning:} the old Bonsai/M1 path demonstrated feasibility, but none of these measurements can satisfy a Qwen prerequisite or gate.
  \end{beamercolorbox}
  \vspace{0.1cm}
  {\scriptsize\color{Red}\textbf{Required next evidence:} a new write-once Qwen/M4 Max preflight receipt and a complete Qwen E0 inside the separate study namespace.}
\end{frame}

\begin{frame}{E1 establishes the prompt-only baseline}
  \small
  \begin{columns}[T,onlytextwidth]
    \begin{column}{0.47\textwidth}
      \textbf{Frozen experiment}\par\vspace{0.15cm}
      One model $\times$ three benchmarks $\times$ three prompts:\par
      P0 neutral; P1 direct answer; P2 calibrated abstention.\par\vspace{0.35cm}
      19,800 deterministic local generations, maximum 48 new tokens, seed 17; 4,800 semantic grades scheduled after generation.
    \end{column}
    \begin{column}{0.47\textwidth}
      \textbf{Why prompt-only comes first}\par\vspace{0.15cm}
      A calibrated instruction may reproduce part of steering's benefit at almost no engineering cost. Without E1, later gains cannot be attributed to internal intervention.\par\vspace{0.35cm}
      E1 also supplies the C/P/I/A labels used to learn probes, construct correct-versus-incorrect activation pairs, and power later paired analyses.
    \end{column}
  \end{columns}
\end{frame}

\begin{frame}{The partial Bonsai E1 cannot be resumed as Qwen evidence}
  \small
  \begin{columns}[T,onlytextwidth]
    \begin{column}{0.62\textwidth}
      \textbf{17,244 of 19,800 pilot generations are retained byte-for-byte}\par\vspace{0.32cm}
      \hbox{%
        \begin{beamercolorbox}[wd=.871\linewidth,ht=2.2ex,dp=0.8ex,center]{progress done}\small 87.1\% superseded pilot\end{beamercolorbox}%
        \begin{beamercolorbox}[wd=.129\linewidth,ht=2.2ex,dp=0.8ex,center]{progress remaining}\small excluded\end{beamercolorbox}}
      \vspace{0.48cm}
      \begin{itemize}\setlength\itemsep{0.12cm}
        \item append-only rows with a resumable chain head
        \item condition-major order across all prompt/benchmark cells
        \item exact checkpoint, split, prompt, runtime, and protocol hashes
        \item exact hashes are recorded in the model-selection amendment
        \item Qwen E1 starts from zero after a new Qwen E0 admission
      \end{itemize}
    \end{column}
    \begin{column}{0.34\textwidth}
      \begin{beamercolorbox}[rounded=true,sep=0.52em]{progress remaining}
        \textbf{Evidence boundary}\par\vspace{0.25cm}
        \footnotesize\textcolor{Green}{\bfseries Can claim:} the abandoned pilot has an immutable, identity-bound record.\par\vspace{0.3cm}
        \textcolor{Red}{\bfseries Cannot claim:} any pilot output is evidence about Qwen or completes E1.
      \end{beamercolorbox}
    \end{column}
  \end{columns}
  \vspace{0.12cm}
  {\tiny\color{Muted}Active namespace: \texttt{artifacts/studies/qwen36-27b-mlx4-m4max48-v1}.}
\end{frame}

\begin{frame}{E2 locates internal factual-outcome separation}
  \small
  \begin{columns}[T,onlytextwidth]
    \begin{column}{0.47\textwidth}
      \textbf{What is implemented}\par\vspace{0.15cm}
      Activation capture at validated sites, compact sharding, calibrated risk probes, layer screening, immutable selection artifacts, and replayable finalization.
    \end{column}
    \begin{column}{0.47\textwidth}
      \textbf{Why E2 exists}\par\vspace{0.15cm}
      Steering should target a location where correct and incorrect outcomes are measurably separable. Choosing a layer first and testing it later would invite cherry-picking.
    \end{column}
  \end{columns}
  \vspace{0.35cm}
  \begin{beamercolorbox}[rounded=true,sep=0.4em]{progress remaining}
    \textbf{Status:} software and replay paths are validated; scientific capture remains blocked until E1 supplies verified outcome labels.
  \end{beamercolorbox}
\end{frame}

\begin{frame}{Static and established baselines come first}
  \small
  \begin{columns}[T,onlytextwidth]
    \begin{column}{0.47\textwidth}
      \textbf{E3 --- static steering}\par\vspace{0.12cm}
      Dense correct-minus-incorrect vectors, layer--alpha--token-scope search, random controls, and label-permutation controls.\par\vspace{0.25cm}
      \textbf{Why:} establish the simplest causal intervention and test whether changes exceed arbitrary perturbation.
    \end{column}
    \begin{column}{0.47\textwidth}
      \textbf{E4 --- established methods}\par\vspace{0.12cm}
      Capability reports and execution receipts for CAA, ITI, ACT, SADI, and TruthX, followed by screening and promotion gates.\par\vspace{0.25cm}
      \textbf{Why:} compare against prior approaches before claiming that the new adaptive controller adds value.
    \end{column}
  \end{columns}
  \vspace{0.3cm}
  {\footnotesize\color{Muted}\textbf{Status:} independently accepted at the software/replay level; experimental execution remains blocked by E1.}
\end{frame}

\begin{frame}{E5 and E6 contain the core causal tests}
  \small
  \begin{columns}[T,onlytextwidth]
    \begin{column}{0.47\textwidth}
      \textcolor{Green}{\bfseries E5 --- adaptive steering: implemented}\par\vspace{0.15cm}
      The matched-budget grid, M1-R/M1-P identity lookup, split provenance, exact alpha binding, signed execution receipts, finalization, and replay checks are implemented.\par\vspace{0.3cm}
      \textbf{Why it matters:} otherwise an adaptive win could come from the wrong baseline, leaked split, or unequal intervention budget.
    \end{column}
    \begin{column}{0.47\textwidth}
      \textcolor{Green}{\bfseries E6 --- knowledge vs. refusal: implemented}\par\vspace{0.15cm}
      Teacher-forced likelihoods, forced-answer ranks, signed row bindings, the registered knowledge-recovery gate, terminal finalizer, and bundle replay are implemented.\par\vspace{0.3cm}
      \textbf{Why it matters:} E6 determines whether E5 changes factual answer selection or only teaches the model to decline.
    \end{column}
  \end{columns}
\end{frame}

\begin{frame}{E7--E10 build the final system}
  \footnotesize
  \renewcommand{\arraystretch}{1.10}
  \begin{tabular}{@{}>{\bfseries\raggedright\arraybackslash}p{0.12\textwidth} >{\raggedright\arraybackslash}p{0.31\textwidth} >{\raggedright\arraybackslash}p{0.47\textwidth}@{}}
    \toprule
    Phase & Planned test & Why it is needed \\
    \midrule
    E7 & Coordinate and sparse-autoencoder steering & Test whether a smaller feature set preserves the factuality effect with fewer collateral changes. \\
    E8 & Protected subspaces for refusal, safety, and language & Prevent a factuality direction from unintentionally changing protected behavior. \\
    E9 & Full prompt $\times$ method factorial & Measure prompt-only, steering-only, combined, and interaction effects under paraphrases. \\
    E10 & Frozen composite policy and untouched confirmation & Combine only independently selected parts, run once, and estimate risk--coverage with utility/safety constraints. \\
    \bottomrule
  \end{tabular}
  \vspace{0.12cm}
  {\scriptsize\color{Muted}Current state: phase orchestration, immutable artifacts, gates, statistics, and replay are implemented; scientific Qwen runs remain.}
\end{frame}

\begin{frame}{Viability is established; efficacy remains open}
  \small
  \begin{columns}[T,onlytextwidth]
    \begin{column}{0.31\textwidth}
      \textcolor{Green}{\bfseries Established}\par\vspace{0.15cm}
      The exact Qwen artifact, runtime policy, study namespace, and E0--E10 software contracts are frozen and replayable.
    \end{column}
    \begin{column}{0.31\textwidth}
      \textcolor{Blue}{\bfseries Next execution}\par\vspace{0.15cm}
      Run live M4 Max preflight, Qwen E0, then a fresh Qwen E1. The preserved Bonsai partial E1 is excluded.
    \end{column}
    \begin{column}{0.31\textwidth}
      \textcolor{Amber}{\bfseries Evidence boundary}\par\vspace{0.15cm}
      Implemented software is not evidence that steering improves factuality; claims require the frozen scientific runs and confirmatory analyses.
    \end{column}
  \end{columns}
  \vspace{0.38cm}
  \begin{beamercolorbox}[rounded=true,sep=0.48em]{progress done}
    \textbf{Honest current conclusion:} the experimental machinery is ready for independent validation and Qwen execution. The central scientific hypotheses remain open.
  \end{beamercolorbox}
\end{frame}

\begin{frame}{Next: complete the evidence chain}
  \footnotesize
  \begin{columns}[T,onlytextwidth]
    \begin{column}{0.47\textwidth}
      \textbf{1. Admit the Qwen runtime}\par
      Download and verify the exact snapshot, run live M4 Max hook preflight, then execute and finalize a fresh Qwen E0.\par\vspace{0.25cm}
      \textbf{2. Run Qwen E1--E4 in order}\par
      Capture only frozen activations, calibrate probes, select static controls, and compare established baselines without touching confirmatory sets.
    \end{column}
    \begin{column}{0.47\textwidth}
      \textbf{3. Execute E5--E8 development}\par
      Fit the adaptive controller, test knowledge recovery, run sparse causal audits, and check protected-behavior non-inferiority.\par\vspace{0.25cm}
      \textbf{4. Reserve E9/E10 only after replay}\par
      Require prerequisite ledgers, registered gates, immutable artifacts, statistical analysis, side-effect checks, and independent validation.
    \end{column}
  \end{columns}
  \vspace{0.2cm}
  \begin{beamercolorbox}[rounded=true,sep=0.3em]{progress done}
    \textbf{Final decision rule:} success means lower hallucination risk at meaningful matched coverage, with correct answers preserved and no material safety, language, or utility regression.
  \end{beamercolorbox}
\end{frame}

\end{document}
